\documentclass[11pt,a4paper]{article}
\usepackage[utf8]{inputenc}
\usepackage[T1]{fontenc}
\usepackage{amsfonts}
\usepackage[french]{babel}
\frenchbsetup{StandardLists=true} % à inclure si on utilise \usepackage[francais]{babel}
\usepackage{enumitem}
\usepackage{amssymb}
\def\nbR{\ensuremath{\mathrm{I\! R}}}
\usepackage{amsmath}
\DeclareMathOperator{\e}{e}

\usepackage[left=2cm,right=2cm,top=2cm,bottom=2cm]{geometry}
\usepackage{multirow}
\usepackage[dvipsnames]{xcolor}
\usepackage{parskip}
\usepackage{caption}




\usepackage[T1]{fontenc}
\usepackage{fouriernc}

\usepackage{graphicx}
\usepackage{setspace}
\singlespacing

\usepackage{tcolorbox}
\usepackage{framed}
\usepackage{multicol}

\usepackage{fancybox}
\usepackage{fancyhdr}
\pagestyle{fancy}
\renewcommand\headrulewidth{1pt}
\fancyhead[L]{Lycée Denis Diderot }
\fancyhead[R]{Classe 1BTS}
\setcounter{page}{1}\fancyfoot[C]{Chapitre 4 - Page \thepage}
\fancyfoot[R]{Année 2025 - 2026}
\fancyfoot[L]{Alexis RUIZ}
\usepackage{pstricks,pst-plot}
\usepackage{tikz}
\usepackage{tkz-tab}
\usetikzlibrary{arrows}
\usepackage{pifont}

\newcommand{\ud}{\mathrm{d}}
\newcommand{\V}{\overrightarrow}
\newcommand{\Rep}{(O;\V{u};\V{v})} 
\newcommand{\cadreblanc}[1]{%
\medskip\framebox[\linewidth][c]{%
\rule{0mm}{#1mm}}\par
\medskip}
\setlength{\columnseprule}{.25mm}
\usepackage{multirow,tabularx,booktabs}
\usepackage{colortbl}
\usepackage{wrapfig}

\usepackage[tikz]{bclogo} 
\usepackage{pgfplots,mathtools}
\usepackage{awesomebox}
\setlength{\aweboxvskip}{0mm}
\newcommand{\defbox}{\awesomebox[black]{2pt}{\faEye}{black}}


\newcommand{\cadre}[1]
{\begin{bclogo}[couleur = white,logo=,  arrondi = 0.3,barre=none]
{%
\rule{0mm}{#1mm}}\par
\medskip
\end{bclogo}}



\begin{document}



\newtcolorbox{mybox}{colback=blue!10!white,
colframe=black!5!black}
\begin{mybox}
\begin{center}    
\textbf{\Large CHAPITRE 4 - DÉRIVATION}\end{center}
\end{mybox}
\vfill
  \begin{bclogo}[couleur = blue!10,logo=\bcinfo,  arrondi = 0.1]
{~Ojectifs~:}
~~
\begin{itemize}

\item Calculer la fonction dérivée d'une fonction.
\item Étudier le sens de variation d'une fonction.
\item Exploiter le tableau de variation d'une fonction
\item Mettre en œuvre un procédé pour déterminer la valeur approchée d'une racine d'une équation 
    
    
\end{itemize}
~~
\end{bclogo}

\vfill


\begin{mybox}
\begin{center}    
\textbf{1. Les dérivées usuelles}\end{center}
\end{mybox}
\vfill



\begin{center}
\renewcommand{\arraystretch}{2.6}
\begin{tabularx}{0.7\linewidth}{|c||*{2}{>{\centering \arraybackslash}X|}}
\hline 
Fonctions & Fonctions dérivées \\ 
\hline 
$a$($a$~constante réelle) &  \\ 
\hline 
$x^n$ &  \\ 
\hline 
$\e^x$ & \\ 
\hline 
$\ln(x)$ &  \\ 
\hline 
$a^x$ ($a>0$)&  \\ 
\hline 
$\sin(x)$ &  \\ 
\hline 
$\cos(x)$ &  \\ 
\hline 
$\tan(x)$ &  \\ 
\hline 
$\arctan(x)$ &  \\ 
\hline 

\end{tabularx}

\end{center}

\vfill
~~

%%%%%%%%%%%%%%%%%%%%%%%%%
\newpage
%%%%%%%%%%%%%%%%%%%%%%%%%

\fcolorbox{black}{orange!30!}{\textbf{{A1.Dérivation de $x^n$ et $a^x$}}}\\[5pt]
Calculer les fonctions dérivées suivantes : \vfill

\begin{center}
\renewcommand{\arraystretch}{3.9}
\begin{tabularx}{0.7\linewidth}{|c||*{2}{>{\centering \arraybackslash}X|}}
\hline 
Fonctions & Fonctions dérivées \\ 
\hline 
$\dfrac{1}{x}$ &  \\ 
\hline 
$\dfrac{1}{x^2}$ &  \\ 
\hline 
$\sqrt{x}$ & \\ 
\hline 
$3^x$ &  \\ 
\hline 
$x^{\frac{3}{2}}$ &  \\ 
\hline 
$x^4$ &  \\ 
\hline 
$4^x$ &  \\ 
\hline 
$3x^2-2x+1$ &  \\ 
\hline 
$\dfrac{x^2}{2}+x+1$ &  \\ 
\hline 

${x^3}+3x^2+x+1$ &  \\ 
\hline 
 & $2x+3$ \\ 
\hline 

\end{tabularx}


\end{center}

\vfill
~~

%%%%%%%%%%%%%%%%%%%%%%%%%%%%%%%%%%%%%%%%
\newpage
%%%%%%%%%%%%%%%%%%%%%%%%%%%%%%%%%%%%%%%%

\begin{mybox}
\begin{center}    
\textbf{2. Les formules de dérivation}\end{center}
\end{mybox}
\vfill

\fcolorbox{black}{blue!10!}{\textbf{2.1. Les formules vues en terminale (ou pas)}}\\
\vfill
Si les fonctions $u$ et $v$ sont deux fonctions dérivables sur le même intervalle $I$. \\[5pt]
 On a les résultats suivants : 
\vfill
\begin{center}
\renewcommand{\arraystretch}{3}
\begin{tabularx}{0.5\linewidth}{|c||*{2}{>{\centering \arraybackslash}X|}}
\hline 
Fonction & Fonction dérivée \\ 
\hline 
$k.u(x)$ &  \\ 
\hline 
$u(x)+v(x)$ &  \\ 
\hline 

$u(x) \times v(x)$ &  \\ 
\hline 
$\dfrac{u(x)}{v(x)}$ &  \\ 
\hline 

\end{tabularx}
\end{center}
\vspace{5mm}
\fcolorbox{black}{green!10!}{\textbf{Exercice 1: Calculer les dérivées des fonctions suivantes }} \\
\begin{multicols}{3}

\begin{center}
{

$f(x)=3x^4-2x^3+6x+34$ 
\vfill
\cadre {100}
$g(x)=\dfrac{x-1}{x+2}$
\vfill
\cadre {100}
$h(x)=(x^2+1)(x-3)$  
\vfill
\cadre {100}



}
\end{center}\end{multicols}

%%%%%%%%%%%%%%%%%%%%%%%%%%%%%%
\newpage
%%%%%%%%%%%%%%%%%%%%%%%%%%%%%%

\fcolorbox{black}{orange!30!}{\textbf{A2. Calculer les dérivées des fonctions suivantes :}} \\
\begin{multicols}{2}

\begin{center}

$f_1(x)=x^4-3x^2+2x-7$ 
\vfill
\cadre{90}
$f_2(x)=3x\ln(x)$
\vfill
\cadre{90}

\end{center}
\end{multicols}

\begin{multicols}{2}
\begin{center}

{
$f_3(x)=(2x-1)(3-5x)$ 
\vfill
\cadre{90}
$f_4(x)= (5-x)\e^x$
\vfill
\cadre{90}
}
\end{center}\end{multicols}

\newpage

\begin{multicols}{2}
\begin{center}

{
$f_5(x)=\dfrac{x-2}{2x+1}$
\vfill
\cadre{90}
$f_6(x)= \dfrac{\ln(x)}{2x}$
\vfill
\cadre{90}
}
\end{center}\end{multicols}

\begin{multicols}{2}
\begin{center}

{
$f_7(x)=\sin(x)\times \cos(x)$
\vfill
\cadre{90}
$f_8(x)= \cos^2x$
\vfill
\cadre{90}
}
\end{center}\end{multicols}

%%%%%%%%%%%%%%%%%%%%%%%%%%%%%%
\newpage
%%%%%%%%%%%%%%%%%%%%%%%%%%%%%%


\fcolorbox{black}{green!10!}{\textbf {Exercice 2 }}\\[5 pt]
Déterminer les dérivées des fonctions suivantes
\begin{multicols}{2}
$f(x)=\e^x(\e^x+3) $
\vfill
\cadre{50}
$g(x)=(\e^x-1)^2$
\vfill
\cadre{50}
\end{multicols}

\noindent
\cadre{35}
\noindent
Ce résultat permet d'obtenir les dérivées des fonctions du tableau 1 en ayant remplacé $x$ par $u(x)$. \\[0.2 cm]
On obtient ainsi les dérivées suivantes : \\

\begin{center}
{
\renewcommand{\arraystretch}{2.5}
\begin{tabularx}{1\linewidth}{|c||p{3.5cm}|||X|}
\hline 
Fonctions & Fonctions dérivées & Exemples \\ 
\hline 
$\left(\e^u\right)$ & & \\ 
\hline 
$u^n$ & & \\ 
\hline 
$\ln u$ & &\\ 
\hline 
$\sin u$ & & \\ 
\hline 
$\cos u $ & & \\ 
\hline 
$\arctan u$ & & \\ 
\hline 
\end{tabularx}
}
\end{center}

\newpage

\fcolorbox{black}{green!10!}{\textbf {Exercice 3 }}\\[5 pt]
Déterminer les dérivées des fonctions suivantes
\begin{multicols}{2}
$f(x)=\e^{3x}(2x+1) $
\vfill
\cadre{50}
$g(x)=(x^2-x+1)\e^{-x}$
\vfill
\cadre{50}
\end{multicols}


\fcolorbox{black}{green!10!}{\textbf {Exercice 4 }}\\[5 pt]
Déterminer les dérivées des fonctions suivantes
\begin{multicols}{2}
$f(x)=\dfrac{1}{4}x^2-\dfrac{1}{4}-\dfrac{1}{2}\ln(x)$
\vfill
\cadre{30}
$g(x)= \ln\left(1+\e^x \right)$
\vfill
\cadre{32}
\end{multicols}



\fcolorbox{black}{green!10!}{\textbf {Exercice 5 }}\\[5 pt]
Déterminer les dérivées des fonctions suivantes
\begin{multicols}{2}
$f(x)=\dfrac{\e^x-\e^{-x}}{\e^x+\e^{-x}}$
\vfill
\cadre{70}
$g(x)= \ln\left(\dfrac{1+x}{1-x}\right)$
\vfill
\cadre{70}
\end{multicols}

\newpage

\fcolorbox{black}{green!10!}{\textbf {Exercice 6 }}\\[5 pt]
Déterminer les dérivées des fonctions suivantes
\begin{multicols}{2}
$f(x)=x^2\ln(x)$
\vfill
\cadre{70}
$g(x)= (x^2-1)\e^{2x}$
\vfill
\cadre{70}

\end{multicols}

\fcolorbox{black}{green!10!}{\textbf {Exercice 7 }}\\[5 pt]
Déterminer les dérivées des fonctions suivantes
\begin{multicols}{2}

$f(x)=\e^{\frac{2x-1}{2-x}}$
\vfill
\cadre{120}
$g(x)= \ln\left(\dfrac{e^x}{e^x+1}\right)$
\vfill
\cadre{120}
\end{multicols}

%%%%%%%%%%%%%%%%%%%%%%%%%%%%%%%%%%%
\newpage           %%%%%%%%%%%%%%%%
%%%%%%%%%%%%%%%%%%%%%%%%%%%%%%%%%%%

\fcolorbox{black}{blue!10!}{\textbf {2.2 Dérivée successives  }}\\[5 pt]
\cadre{55}
\cadre{15}

\fcolorbox{black}{orange!30!}{\textbf {A3. Déterminer $f^{(5)}$ avec $f(x)=\cos(3x)$}}\\[5 pt]
\cadre{35}

\fcolorbox{black}{blue!10!}{\textbf {3. Les principales utilisations de la dérivée  }}\\[5 pt]
\fcolorbox{black}{blue!10!}{\textbf {3.1 Équation des tangentes et variations d'une fonction  }}\\[5 pt]
\cadreblanc{65}

%%%%%%%%%%%%%%%%%%%%%%%%%%%%%%%
\newpage
%%%%%%%%%%%%%%%%%%%%%%%%%%%%%%%

\fcolorbox{black}{orange!30!}{\textbf {A4. Une équation de la tangente}}\\[5 pt]
Soit $f$ une fonction définie sur $\mathbb{R}$ par $f(x)=(2x+1)\e^x$ et C sa courbe représentative. 
Donner une équation de la tangente à C au point d'abscisse $x=0$.Vérifier votre résultat à l'aide de la calculatrice en traçant la courbe représentative de $f$ et la tangente trouvée. \\[5 pt]
\cadre{60}

\vfill

\fcolorbox{black}{blue!10!}{\textbf {3.2 Théorèmes relatifs au sens de variation d'une fonction }}\\[5 pt]
\defbox{\fbox{
\begin{minipage}{14cm}{

Si sur un intervalle $I$ la dérivée $f'$ de la fonction $f$ est positive alors la fonction $f$ est croissante sur $I$.\\[5 pt]
Si sur un intervalle $I$ la dérivée $f'$ de la fonction $f$ est négative alors la fonction $f$ est décroissante sur $I$.
}

\end{minipage}
}}

\vfill 

\fcolorbox{black}{green!10!}{\textbf {Exercice 8 - Faire le point avec un QCM}}\\[5 pt]
Soit $f$ la fonction définie sur $\mathbb{R}$ par $f(x)=\e^{2x}(x^2+x-1)$ et $C$ sa courbe représentative : \\
\begin{multicols}{2}

1. Dérivée 
\begin{itemize}[label=$\square$]
\item   a) $f'(x)=(x^2+3x)\e^{2x}$
\item   b) $f'(x)=(2x^2+4x-1)\e^{2x}$
\item   c) $f'(x)=(2x+1)\e^{2x}$
\end{itemize}
\begin{center}
\noindent\rule{5cm}{0.4pt}
\end{center}


2. Une équation de la droite tangente à $C$ au point d'abscisse $a=0$ est :
\begin{itemize}[label=$\square$]
\item   a) $y=3x+1$
\item   b) $y=-1$
\item   c) $y=-x-1$
\end{itemize}


3. $f$ est strictement croissante sur :
\begin{itemize}[label=$\square$]
\item   a) $\mathbb{R}$
\item   b) $[1;+\infty[$
\item   c) $[0;1]$
\end{itemize}

\begin{center}
\noindent\rule{5cm}{0.4pt}
\end{center}

4. $f$ est strictement positive sur :  
\begin{itemize}[label=$\square$]
\item   a) $\mathbb{R}$
\item   b) $[1;+\infty[$
\item   c) $[0;1]$
\end{itemize}

\end{multicols}

\vfill 

\newpage 
%%%%%%%%%%%%%%%%%%%%%%%%%%%%%%%%%%%%%%%


\fcolorbox{black}{green!10!}{ \textbf {Exercice 9 - Avec plusieurs fonctions dérivées}}\\
\normalsize

\flushleft{Soit la fonction $f$ définie sur $\mathbf{R}$ par $f(x)=(3x-1)\e^{2x}$ : \\
1. Calculer $f'(x)$ et $f''(x).$
\begin{multicols}{2}
\cadre{50}
\cadre{50}
\end{multicols}

2.Montrer que pour tout réel $x$, $f''(x) -f'(x) - 2f(x) = 9\e^{2x}.$
\cadreblanc{50}

\fcolorbox{black}{green!10!}{\textbf {Exercice 10 - Équation de la tangente}}\\
Écrire une équation de la tangente à la courbe représentative de la fonction $f$ au point d'abscisse $a$ donné : 

\begin{multicols}{2}
$f(x)= \ln x + 2x - 3$ en $a=1$
\cadre {70}
$f(x)=\e^x(3-2x)$ en $a=0$
\cadre {70}
\end{multicols}


%%%%%%%%%%%%%%%%%%%%%%%%%%%%%%%%%%%%%%
\newpage
%%%%%%%%%%%%%%%%%%%%%%%%%%%%%%%%%%%%%%

    
 \fcolorbox{black}{blue!10!}{\textbf {3.3. Le théorème des valeurs intermédiaires.   }}\\[0.3cm]
 
\cadre {25}

\cadre {25}
\textbf {Exemple  }\\[0.1cm]

Étudier le sens de variation de la fonction $f$définie sur $\mathbf{R}$ par $f(x) = \e^{-x}(x+2)$ et en déduire sans 
calcul que l'équation $f(x)=0$ n'admet qu'une solution dans $\mathbf{R}$.


\cadre {140}

%%%%%%%%%%%%%%%%%%%%%%%%%%%%%%%%%%%%%%%%%%
\newpage
%%%%%%%%%%%%%%%%%%%%%%%%%%%%%%%%%%%%%%%%%%


\fcolorbox{black}{green!10!}{\textbf {Exercice 11 :} }\\ [5 pt]

Soit la fonction $f$ définie sur $\mathbf{R}$ par $f(x)=\e^x(4-2x)$

On nomme C la courbe représentative dans un repère orthonormé.


1. Déterminer les limites de $f$ en $-\infty$ et  $+\infty$.
\cadre{40}


2. Calculer $f'(x)$ et étudier son signe sur $\mathbf{R}$
\cadre{30}
\cadre{30}

3. Donner le tableau de variation de $f$.
\cadre{30}

4. Écrire une équation de la tangente à C au point d'abscisse $a=1$
\cadre{30}


%%%%%%%%%%%%%%%%%%%%%%%%%%%%%%%%%%
\newpage
%%%%%%%%%%%%%%%%%%%%%%%%%%%%%%%%%%


\fcolorbox{black}{green!10!}{\textbf {Exercice 12 :} }\\ [5 pt]

Soit la fonction $f$ définie sur $\mathbf{R}$ par $f(x)=\e^{-x}(3-4x)$

On nomme C la courbe représentative dans un repère orthonormé.


1. Déterminer les limites de $f$ en $-\infty$ et  $+\infty$.
\cadre{40}


2. Calculer $f'(x)$ et étudier son signe sur $\mathbf{R}$
\cadre{30}
\cadre{30}

3. Donner le tableau de variation de $f$.
\cadre{30}

4. Écrire une équation de la tangente à C au point d'abscisse $a=0$
\cadre{30}



\end{document}

